\subsection{Variational principles for the spectrum of positive operators}

In this number, E is a nonzero complex Hilbert space.

PROPOSITION 15. --- Let u be a self-adjoint partial operator on E. Suppose that u is bounded below by \(c \in \mathbf{R}\) (def. 7 of IV, p. 237). We have

\[
\inf (\mathrm{Sp} (u)) = \inf _ {x \in \mathrm{dom} (u) - \{0 \}} \frac {\langle x \mid u (x) \rangle}{\| x \| ^ {2}} \in [ c, + \infty [.
\]

Let m be the right-hand side of the equality to be proved. Let \(x \in E\) and \(\mu_{x}\) be the spectral measure of x relative to u. We have

\[
\begin{array}{r l} \langle x \mid u (x) \rangle & = \int_ {\mathrm{Sp} (u)} t d \mu_ {x} (t) \\ & \geqslant \inf (\mathrm{Sp} (u)) \int_ {\mathrm{Sp} (u)} d \mu_ {x} (t) = \inf (\mathrm{Sp} (u)) \| x \| ^ {2}, \end{array}
\]

(formula (6), p. 271), hence \(\inf(\mathrm{Sp}(u)) \leqslant m\) . Conversely, the partial operator \(u - m\) is positive, therefore \(\inf(\mathrm{Sp}(u)) - m = \inf(\mathrm{Sp}(u - m \cdot 1_{\mathrm{E}})) \geqslant 0\) (prop. 17 of IV, p. 248).

Proposition 9 of I, p. 139 corresponds to the particular case of this proposition when u is a Hermitian element of \(\mathcal{L}(\mathrm{E})\) , which is then necessarily bounded below. In this case, we also have

\[
\sup (\operatorname{Sp} (u)) = \sup _ {x \in \operatorname{dom} (u) - \{0 \}} \frac {\langle x \mid u (x) \rangle}{\| x \| ^ {2}}
\]

(loc. cit.); if u does not belong to \(\mathcal{L}(\mathrm{E})\) then this supremum is \(+\infty\) .

DEFINITION 10. --- Let u be a normal partial operator on E. A complex number \(\lambda \in \mathrm{Sp}(u)\) belongs to the sensitive spectrum of \(u\) if \(\lambda\) is isolated in \(\mathrm{Sp}(u)\) and if \(\lambda\) is an eigenvalue of finite multiplicity.

We denote by \(\mathrm{Sp}_{\mathrm{s}}(u)\) the sensitive spectrum of u. Its complement in \(\mathrm{Sp}(u)\) is called the essential spectrum of u and denoted \(\mathrm{Sp}_{\mathrm{e}}(u)\) .

The essential spectrum of a normal partial operator u on E is closed in C, since \(\mathrm{Sp}(u)\) is closed in C and the complement of the essential spectrum is open in \(\mathrm{Sp}(u)\). The sensitive spectrum of u is not necessarily closed in C (Exercise 36 of IV, p. 369).

Lemma 9. --- Let E be a complex Hilbert space and let u be a normal partial operator on E. Let \(\lambda \in \mathrm{Sp}(u)\) . We have \(\lambda \in \mathrm{Sp}_{\mathrm{s}}(u)\) if and only if there exists an open neighborhood V of \(\lambda\) in C such that the spectral projector \(p_{\mathrm{V}}\) of u defined by V has finite rank.

Let \(\lambda \in \mathrm{Sp}_{\mathrm{s}}(u)\). There exists an open neighborhood \(V\) of \(\lambda\) in \(\mathbf{C}\) such that \(\mathrm{Sp}(u) \cap V = \{\lambda\}\) ; the spectral projector \(p_V = p_{\{\lambda\}}\) is then of finite rank (cor. of prop. 9 of IV, p. 278).

Conversely, suppose that there exists an open neighborhood V of \(\lambda\) such that the spectral projector \(p_{V}\) has finite rank \(n \in N\). By the corollary of Proposition 10 of IV, p. 279, the intersection \(V \cap \mathrm{Sp}(u)\) contains at most n elements, therefore \(\lambda\) is isolated in \(\mathrm{Sp}(u)\). This is therefore an eigenvalue of u of spectral multiplicity at most n, by the corollary of prop. 9 of IV, p. 278.

In the rest of this number, we assume that E is infinite-dimensional; the analogues of the results below when E is finite-dimensional are deduced from No. 4 of IV, p. 153.

Let u be a self-adjoint operator bounded below on E. Let c be a lower bound of u; the spectrum of u is contained in \([c, +\infty[\) (prop. 15). Suppose that the essential spectrum of u is not empty. We then denote by \(\varrho_{e}\) the lower bound of the essential spectrum of u, so that \(\varrho_{e} \geqslant c\) and \(\varrho_{e}\) is an element of the spectrum of \(u\) . We denote \(\mathrm{Sp}_{\mathrm{h}}(u) = \mathrm{Sp}(u) \cap [\varrho_{\mathrm{e}}, +\infty[\); it is a closed subset of \(\mathrm{Sp}(u)\) , hence also of \(\mathbf{C}\) , such that \(\inf (\mathrm{Sp}_{\mathrm{h}}(u)) = \varrho_{\mathrm{e}}\) ; it is called the upper part of the spectrum of \(u\) . If the essential spectrum of \(u\) is empty, we denote \(\mathrm{Sp}_{\mathrm{h}}(u) = \varnothing\) and \(\varrho_{\mathrm{e}} = +\infty\) .

The intersection \(\mathrm{Sp}_{\mathrm{b}}(u) = \mathrm{Sp}(u) \cap [0, \varrho_{\mathrm{e}}[\) is contained in the sensitive spectrum of \(u\) and its elements are therefore isolated eigenvalues of finite multiplicity; this is called the lower part of the spectrum of \(u\) . Let \(\mathrm{E}_{\mathrm{b}}\) be the closed subspace of E generated by the eigenspaces corresponding to the eigenvalues \(\lambda \in \mathrm{Sp}_{\mathrm{b}}(u)\) . By the cor. of prop. 9 of IV, p. 278 and prop. 8 of IV, p. 278, the space \(\mathrm{E}_{\mathrm{b}}\) is the image of the spectral projector \(p_{\mathrm{Sp}_{\mathrm{b}}(u)}\) defined by \(\mathrm{Sp}_{\mathrm{b}}(u)\) . Since \(\mathrm{Sp}(u)\) is the disjoint union of \(\mathrm{Sp}_{\mathrm{b}}(u)\) and \(\mathrm{Sp}_{\mathrm{h}}(u)\) , the orthogonal complement \(\mathrm{E}_{\mathrm{h}}\) of \(\mathrm{E}_{\mathrm{b}}\) is the image of the spectral projector \(p_{\mathrm{Sp}_{\mathrm{h}}(u)}\) defined by \(\mathrm{Sp}_{\mathrm{h}}(u)\) . If \(\mathrm{Sp}_{\mathrm{b}}(u)\) is finite, then the space \(\mathrm{E}_{\mathrm{b}}\) is finite-dimensional; the space \(\mathrm{E}_{\mathrm{h}}\) is then nonzero, since E is assumed to be infinite-dimensional.

Lemma 10. --- One has

\[
\langle x \mid u (x) \rangle \geqslant \varrho_ {\mathrm{e}} \| x \| ^ {2}.\tag{18}
\]

for all \(x \in \mathrm{E}_{\mathrm{h}} \cap \operatorname{dom}(u)\) .

If \(x \in \mathrm{E}_{\mathrm{h}} \cap \mathrm{dom}(u)\) , then the support of the spectral measure \(\mu_{x}\) of \(x\) relative to \(u\) is contained in the interval \([\varrho_{\mathrm{e}}, +\infty[\) (prop. 9, \(a\) ) of IV, p. 278), hence

\[
\langle x \mid u (x) \rangle = \int_ {\mathbf {R}} t d \mu_ {x} (t) \geqslant \varrho_ {\mathrm{e}} \| x \| ^ {2}.
\]

Lemma 11. --- Suppose that \(E_{b}\) is finite-dimensional. Then, for every real number \(\varepsilon > 0\), the image of the spectral projector of u defined by \([\varrho_{e}, \varrho_{e} + \varepsilon]\) is infinite-dimensional.

If \(E_{b}\) is finite-dimensional, then the essential spectrum of u is nonempty, hence \(\varrho_{e}\) is finite and belongs to \(\mathrm{Sp}_{\mathrm{e}}(u)\) . Moreover, the lower spectrum \(\mathrm{Sp}_{\mathrm{b}}(u)\) is finite, so there exists \(\delta > 0\) such that \([\varrho_{e} - \delta, \varrho_{e} + \delta] \cap \mathrm{Sp}(u) \subset [\varrho_{e}, \varrho_{e} + \delta]\) . The assertion then follows from Lemma 9.

PROPOSITION 16. --- The set \(\mathrm{Sp}_{\mathrm{b}}(u) \subset [c, \varrho_{\mathrm{e}}[\) is countable, discrete, and closed in \([0, \varrho_{\mathrm{e}}[\) . It is the set of values of a strictly increasing sequence \((\nu_{k})_{0 \leqslant k < \mathrm{Card}(\mathrm{Sp}_{\mathrm{b}}(u))}\) of positive real numbers; if \(\mathrm{Sp}_{\mathrm{b}}(u)\) is infinite, then the sequence \((\nu_{k})\) converges to \(\varrho_{\mathrm{e}}\) .

Let \(\mathrm{T} = \mathrm{Sp}_{\mathrm{b}}(u) \cap [0, \varrho_{\mathrm{e}}[\) . It is a closed and discrete subset of \([c, \varrho_{\mathrm{e}}[\) by definition. For every integer \(i \geqslant 1\) , the set \(\mathrm{Sp}_{\mathrm{b}}(u) \cap [c, \varrho_{\mathrm{e}} - 1/i]\) is compact and discrete, hence finite. Since T is the union of these sets for \(i \geqslant 1\) , the set T is countable.

This concludes the proof if \(\mathrm{Sp}_{\mathrm{b}}(u)\) is finite. If \(\mathrm{Sp}_{\mathrm{b}}(u)\) is infinite, one concludes by applying Lemma 2 of III, p. 90 to the image S of T under the map \(\lambda \mapsto \varrho_{e} - \lambda\) .

For every integer \(k\) such that \(0 \leqslant k < \text{Card}(\text{Sp}_b(u))\) , we denote by \(n_k \geqslant 1\) the multiplicity of the eigenvalue \(\nu_k\) of \(u\) . We denote \(\overline{\mathbf{N}} = \mathbf{N} \cup \{+\infty\} \subset \overline{\mathbf{R}}\) . Let \(\mathbf{M} \in \overline{\mathbf{N}}\) be the sum of the multiplicities \(n_k\) . This is the Hilbertian dimension of \(\mathbf{E}_b\) , if one agrees to say that a space of infinite Hilbertian dimension has dimension \(+\infty \in \overline{\mathbf{N}}\) .

For every integer n such that \(0 \leqslant n < M\) , one defines \(\lambda_{n}(u) = \nu_{k}\) , where \(k \geqslant 0\) is the unique integer such that

\[
n _ {0} + \dots + n _ {k - 1} \leqslant n <   n _ {0} + \dots + n _ {k}.
\]

If \(n \in \mathbf{N}\) satisfies \(n \geqslant \mathrm{M}\), we set \(\lambda_n(u) = \varrho_{\mathrm{e}}\) . This case can occur only if \(\mathrm{Sp}_{\mathrm{b}}(u)\) is finite, in which case \(\varrho_{\mathrm{e}}\) is finite, since E is assumed to be of infinite dimension.

By construction, the sequence \((\lambda_{n}(u))_{n\in\mathbf{N}}\) is increasing and, for every element \(\lambda\) of \(\mathrm{Sp}_{\mathrm{b}}(u)\) , the number of integers n such that \(\lambda_{n}(u)=\lambda\) is equal to the multiplicity of \(\lambda\) as an eigenvalue of u. The sequence \((\lambda_{n}(u))_{n\in\mathbf{N}}\) tends to \(+\infty\) if and only if the essential spectrum of u is empty.

For every \(n \in N\) , we denote by \(F_{n}\) (resp. \(F_{n}^{u}\) ) the set of vector subspaces \(F \subset E\) of dimension n (resp. the set of vector subspaces \(F \subset \text{dom}(u)\) of dimension n).

Let \(n \in \mathbf{N}\) such that \(n < \mathrm{M}\) and \(\mathrm{F} \in \mathcal{F}_n^u\) . One says that \(\mathrm{F}\) is adapted to \(u\) if \(\mathrm{F}\) admits an orthonormal basis \((f_i)_{0 \leqslant i \leqslant n-1}\) such that \(u(f_i) = \lambda_i(u)f_i\) for \(0 \leqslant i \leqslant n-1\) .

Let \(F \in F_{n}^{u}\) adapted to u. The space F is contained in \(E_{b}\) ; it is generated by eigenvectors of u for eigenvalues \(\lambda \leqslant \lambda_{n-1}(u)\) , and it contains the eigenspaces corresponding to the eigenvalues \(\lambda < \lambda_{n-1}(u)\) . Moreover, for every universally measurable subset A of \(\mathrm{Sp}(u)\) , the space F is stable under the spectral projector \(p_{A}\) defined by A (prop. 9, c) of IV, p. 278).

Lemma 12. --- Let \(F \in F_{n}^{u}\) be a subspace adapted to u. Every eigenvector of u belonging to the space \(F^{\circ} \cap E_{b}\) has an eigenvalue \(\geqslant \lambda_{n}(u)\) .

Let \(\ell\) be such that \(\lambda_{n-1}(u) = \nu_{\ell}\) . Let \(x \in \mathrm{F}^{\circ} \cap \mathrm{E}_{\mathrm{b}}\) be an eigenvector of \(u\) for an eigenvalue \(\lambda\) and let \(k < \operatorname{Card}(\operatorname{Sp}_{\mathrm{b}}(u))\) be such that \(\lambda = \nu_{k}\) .

The condition \(k < \ell\) is impossible, because \(\nu_k < \nu_\ell\) , and the space F would then contain the eigenspace for the eigenvalue \(\nu_k\) , contradicting the fact that \(x\) is orthogonal to F. If \(k = \ell\) , then \(x\) is an eigenvector for the eigenvalue \(\lambda_{n-1}(u)\) ; since \(x \in \mathrm{F}^\circ\) , the multiplicity \(n_k\) is strictly greater than the number of integers \(i < n\) such that \(\lambda_i(u) = \nu_k\) , which implies that \(\lambda_n(u) = \lambda_{n-1}(u)\) . Finally, if \(k > \ell\) , we have \(\lambda = \nu_k > \nu_\ell = \lambda_{n-1}(u)\) , hence \(\lambda \geqslant \lambda_n(u)\) .

For every subspace F of E, we denote by

\[
\widetilde{r}_{\mathrm{F}}(u) = \inf_{\substack{x\in \operatorname{dom}(u)\cap \mathrm{F}^{\circ}\\ x\neq 0}}\frac{\langle x\mid u(x)\rangle}{\|x\|^{2}},
\]

\[
\widetilde{\mathrm{R}}_{\mathrm{F}}(u) = \sup_{\substack{x\in \operatorname{dom}(u)\cap \mathrm{F}\\ x\neq 0}}\frac{\langle x\mid u(x)\rangle}{\|x\|^ {2}}.
\]

PROPOSITION 17. --- a) For every integer \(n \in \mathbf{N}\) , one has

\[
\lambda_ {n} (u) = \sup _ {\mathrm{F} \in \mathcal {F} _ {n}} \widetilde {r} _ {\mathrm{F}} (u) = \inf _ {\mathrm{F} \in \mathcal {F} _ {n + 1} ^ {u}} \widetilde {\mathrm{R}} _ {\mathrm{F}} (u);
\]

\begin{enumerate}
\def\labelenumi{\alph{enumi})}
\setcounter{enumi}{1}
\item
  For every integer \(n < \mathrm{M}\) and for every subspace \(\mathrm{F} \in \mathcal{F}_n^u\) adapted to \(u\) , we have \(\lambda_n(u) = \widetilde{r}_{\mathrm{F}}(u)\) ;
\item
  For every integer \(n < \mathrm{M}\) and for every subspace \(\mathrm{F} \in \mathcal{F}_{n+1}^{u}\) adapted to \(u\) , we have \(\lambda_{n}(u) = \widetilde{\mathrm{R}}_{\mathrm{F}}(u)\) .
\end{enumerate}

There exists an orthonormal basis \((e_{n})_{0\leqslant n<M}\) of the space \(E_{b}\) such that \(e_{n}\in\operatorname{dom}(u)\) and \(u(e_{n})=\lambda_{n}(u)e_{n}\) for every n such that \(0\leqslant n<M\) . For every \(x\in E_{b}\cap\operatorname{dom}(u)\) , we therefore have

\[
\langle x \mid u (x) \rangle = \sum_ {0 \leqslant n <   M} \lambda_ {n} (u) | \langle e _ {n} \mid x \rangle | ^ {2}.
\]

For every integer \(n\) such that \(1 \leqslant n < M + 1\) , let \(F_n\) be the subspace of dimension \(n\) of \(E_b\) generated by \((e_0, \ldots, e_{n-1})\) . We have \(F_n \subset \text{dom}(u)\) and \(F_n\) is adapted to \(u\) .

Let \(n \in \mathbf{N}\) and \(\mathrm{F} \in \mathcal{F}_n\) . Let us prove that \(\widetilde{r}_{\mathrm{F}}(u) \leqslant \lambda_n(u)\) and, consequently, that

\[
\sup _ {\mathrm{F} \in \mathcal {F} _ {n}} \widetilde {r} _ {\mathrm{F}} (u) \leqslant \lambda_ {n} (u).\tag{19}
\]

If \(0 \leqslant n < M\) , in particular if M is infinite, then the restriction to \(\mathrm{F}_{n+1}\) of the orthogonal projection of E onto F is not injective, hence there exists \(x \neq 0\) in \(F_{n+1}\) orthogonal to F. Since

\[
\langle x \mid u (x) \rangle = \sum_ {0 \leqslant i \leqslant n} \lambda_ {i} (u) | \langle e _ {i} \mid x \rangle | ^ {2} = \lambda_ {n} (u) | \langle e _ {n} \mid x \rangle | ^ {2} \leqslant \lambda_ {n} (u) \| x \| ^ {2},
\]

we have \(\widetilde{r}_{\mathrm{F}}(u) \leqslant \lambda_n(u)\) .

If \(M \in N\) and \(n \geqslant M\) , then for every real number \(\varepsilon > 0\) , there exists x of norm 1 in \(\text{dom}(u)\) which is orthogonal to F and whose spectral measure \(\mu_{x}\) has support contained in \([\varrho_{e}, \varrho_{e} + \varepsilon]\) (lemma 11 and prop. 9, a) of IV, p. 278), whence

\[
\widetilde {r} _ {\mathrm{F}} (u) \leqslant \langle x \mid u (x) \rangle = \int_ {\operatorname{Sp} (u)} t d \mu_ {x} (t) \leqslant \varrho_ {\mathrm{e}} + \varepsilon = \lambda_ {n} (u) + \varepsilon
\]

by definition. Since \(\varepsilon > 0\) is arbitrary, we therefore have \(\widetilde{r}_{\mathrm{F}}(u) \leqslant \lambda_n(u)\) . Inequality (19) is therefore established.

Let \(n \in \mathbf{N}\) and \(\mathrm{F} \in \mathcal{F}_{n+1}^{u}\) . Let us prove that \(\widetilde{\mathrm{R}}_{\mathrm{F}}(u) \geqslant \lambda_{n}(u)\) and, consequently, that

\[
\inf _ {\mathrm{F} \in \mathcal {F} _ {n + 1} ^ {u}} \widetilde {\mathrm{R}} _ {\mathrm{F}} (u) \geqslant \lambda_ {n} (u).\tag{20}
\]

If \(0 \leqslant n < M\) , observe that the restriction to F of the orthogonal projector onto \(\mathrm{F}_n\) is not injective, hence there exists a vector \(x \neq 0\) in F orthogonal to \(\mathrm{F}_n\) . Set \(x_{\mathrm{b}} = p_{\mathrm{Sp}_{\mathrm{b}}(u)}(x)\) and \(x_{\mathrm{h}} = p_{\mathrm{Sp}_{\mathrm{h}}(u)}(x)\) . We therefore have \(x = x_{\mathrm{b}} + x_{\mathrm{h}}\) . The elements \(x_{\mathrm{b}}\) and \(x_{\mathrm{h}}\) belong to the domain of \(u\) (prop. 9, c) of IV, p. 278) and are orthogonal to \(\mathrm{F}_n\) . We have the lower bound

\[
\langle x _ {\mathrm{b}} \mid u (x _ {\mathrm{b}}) \rangle = \sum_ {i \geqslant n} \lambda_ {i} (u) | \langle e _ {i} \mid x _ {\mathrm{b}} \rangle | ^ {2} \geqslant \lambda_ {n} (u) \| x _ {\mathrm{b}} \| ^ {2},
\]

and \(\langle x_{\mathrm{h}}|u(x_{\mathrm{h}})\rangle \geqslant \varrho_{\mathrm{e}}\| x_{\mathrm{h}}\|^2\) (formula (18)), whence

\[
\begin{array}{r l} & {\langle x \mid u (x) \rangle = \langle x _ {\mathrm{b}} \mid u (x _ {\mathrm{b}}) \rangle + \langle x _ {\mathrm{h}} \mid u (x _ {\mathrm{h}}) \rangle} \\ & {\qquad \geqslant \lambda_ {n} (u) \| x _ {\mathrm{b}} \| ^ {2} + \varrho_ {\mathrm{e}} \| x _ {\mathrm{h}} \| ^ {2} \geqslant \lambda_ {n} (u) \| x \| ^ {2}.} \end{array}
\]

If M is finite and \(n \geqslant M\) , there exists \(x \neq 0\) in F orthogonal to \(E_{b}\) , hence \(x \in E_{h}\) , and

\[
\langle x \mid u (x) \rangle \geqslant \varrho_ {\mathrm{e}} \| x \| ^ {2} = \lambda_ {n} (u) \| x \| ^ {2}
\]

by (18). Inequality (20) is therefore proved.

We shall now prove the assertions \(b\) ) and \(c\) ), which imply that inequalities (19) and (20) are equalities when \(n < M\) .

Let us prove assertion \(b\) ). Let \(\mathrm{F} \in \mathcal{F}_n^u\) be a space adapted to \(u\) . We have the Hilbert sum

\[
\mathrm{F}^{\circ} = \mathrm{E}_{\mathrm{h}}\oplus \bigoplus_{\substack{\lambda \in \operatorname{Sp}_{\mathrm{b}}(u)\\ \lambda >\lambda_{n - 1}(u)}}\operatorname{Ker}(u - \lambda \cdot 1_{\mathrm{E}})\oplus (\mathrm{F}^{\circ}\cap \operatorname{Ker}(u - \lambda_{n - 1}(u)\cdot 1_{\mathrm{E}})).
\]

Let \(x \in \mathrm{F}^{\circ} - \{0\}\) . Let us write \(x = x_{\mathrm{h}} + y + z\) , where \(x_{\mathrm{h}} \in \mathrm{E}_{\mathrm{h}}\) and \(y\) (resp. \(z\) ) belongs to the second (resp. third) space of the above decomposition. Using (18) again and the fact that every eigenvalue \(\lambda > \lambda_{n-1}(u)\) of \(u\) is \(\geqslant \lambda_n(u)\) , we obtain

\[
\begin{array}{r l} & {\langle x \mid u (x) \rangle = \langle x _ {\mathrm{h}} \mid u (x _ {\mathrm{h}}) \rangle + \langle y \mid u (y) \rangle + \langle z \mid u (z) \rangle} \\ & {\qquad \geqslant \varrho_ {\mathrm{e}} \| x _ {\mathrm{h}} \| ^ {2} + \lambda_ {n} (u) \| y \| ^ {2} + \lambda_ {n - 1} (u) \| z \| ^ {2}.} \end{array}
\]

If \(z \neq 0\) , then by Lemma 12, we have \(\lambda_n(u) = \lambda_{n-1}(u)\) . We therefore deduce that \(\langle x | u(x) \rangle \geqslant \lambda_n(u) \| x \|^2\) , whence \(\widetilde{r}_{\mathrm{F}}(u) \geqslant \lambda_n(u)\) . Combined with (19), this implies the assertion \(b\) ).

Let us prove assertion \(c\) ). Let \(\mathrm{F} \in \mathcal{F}_{n+1}^{u}\) be a subspace adapted to \(u\) . Let \((f_{i})_{0 \leqslant i \leqslant n}\) be an orthonormal basis of \(\mathrm{F}\) such that \(u(f_{i}) = \lambda_{i}(u)f_{i}\) for \(0 \leqslant i \leqslant n\) . For every \(x \in \mathrm{F}\) , we have

\[
\langle x \mid u (x) \rangle = \sum_ {0 \leqslant i \leqslant n} \lambda_ {i} (u) | \langle f _ {i} \mid x \rangle | ^ {2} \leqslant \lambda_ {n} (u) \| x \| ^ {2},
\]

with equality if \(x = f_{n}\) , hence \(\widetilde{\mathrm{R}}_{\mathrm{F}}(u) = \lambda_{n}(u)\) .

Let us finally prove that (19) and (20) are equalities when \(n \geqslant M\) . In this case, M is finite, hence \(E_{b}\) is contained in the domain of u; moreover, we have \(\lambda_{n}(u) = \varrho_{\mathrm{e}}\) by definition.

There exists \(F \in F_{n}\) containing \(E_{b}\) . Every element \(x \neq 0\) of \(\text{dom}(u)\) orthogonal to F is therefore orthogonal to \(E_{b}\) , and satisfies

\[
\frac {\langle x \mid u (x) \rangle}{\| x \| ^ {2}} \geqslant \varrho_ {\mathrm{e}}
\]

(formula (18)), hence \(\widetilde{r}_{\mathrm{F}}(u) \geqslant \varrho_{\mathrm{e}} = \lambda_n(u)\) , and consequently

\[
\sup _ {\mathrm{F} \in \mathcal {F} _ {n}} \widetilde {r} _ {\mathrm{F}} (u) \geqslant \lambda_ {n} (u).
\]

Let \(\varepsilon > 0\) . Since \(\mathrm{E_b}\) is finite-dimensional, there exists by Lemma 11 an orthonormal family \((x_i)_{i \in \mathrm{I}}\) of elements of E such that the subspace F generated by \(\mathrm{E_b}\) and \((x_i)_{i \in \mathrm{I}}\) has dimension \(n + 1\) and is contained in \(\operatorname{dom}(u)\) , and such that \(\langle x_{i} \mid u(x_{i})\rangle \leqslant \varrho_{\mathrm{e}} + \varepsilon\) for all \(i \in I\) . We then have \(\widetilde{\mathrm{R}}_{\mathrm{F}}(u) \leqslant \varrho_{\mathrm{e}} + \varepsilon\) . Since \(\varepsilon > 0\) is arbitrary, we conclude that

\[
\inf _ {\mathrm{F} \in \mathcal {F} _ {n + 1} ^ {u}} \widetilde {\mathrm{R}} _ {\mathrm{F}} (u) \leqslant \varrho_ {\mathrm{e}} = \lambda_ {n} (u).
\]

The proposition is proved.

\subsection{Compact perturbation and essential spectrum}

In this number, E is an infinite-dimensional complex Hilbert space.

Lemma 13. --- Let I be an orthonormal family in E. The family I converges weakly to 0 along the filter of complements of finite subsets of I.

Let \(x \in E\) . By Bessel's inequality (EVT, V, p. 21, prop. 4) and TG, IV, p. 37, th. 1, the family \((|\langle e_{i} | x\rangle|^{2})_{i\in I}\) is summable, hence \(\langle e_{i} | x\rangle\) tends to 0 along the filter of complements of finite subsets of I (TG, III, p. 38, prop. 1).

PROPOSITION 18. --- Let \(u\) be a self-adjoint partial operator on E and let \(\lambda\) be a real number. We have \(\lambda \in \mathrm{Sp}_{\mathrm{e}}(u)\) if and only if there exists an orthonormal sequence \((x_{n})_{n \in \mathbf{N}}\) in E such that \(u(x_{n}) - \lambda x_{n}\) tends to 0 in E.

Suppose first that \(\lambda \in \mathrm{Sp}_{\mathrm{e}}(u)\) . If the spectral projector of \(u\) relative to \(\{\lambda\}\) is of infinite rank, every orthonormal sequence \((x_{n})\) in its image answers the question (cf. cor. of prop. 9 of IV, p. 278). In what follows we shall assume that this is not the case.

For every \(k \in \mathbf{N}\) , let \(\mathrm{J}_k\) be the set of \(t \in [\lambda - 1, \lambda + 1]\) such that

\[
{\frac {1}{k + 2}} <   | t - \lambda | \leqslant {\frac {1}{k + 1}}.
\]

The sets \(J_{k}\) are pairwise disjoint. Moreover, for every integer \(K \in N\) , the sets \((\mathrm{J}_{k})_{k \geqslant K}\) form a partition of the set

\[
\mathrm{I} _ {\mathrm{K}} = [ \lambda - 1 / (\mathrm{K} + 1), \lambda + 1 / (\mathrm{K} + 1) ] - \{0 \}.
\]

Since the spectral projector of \(u\) relative to \(\mathrm{I_K} \cup \{0\}\) has infinite rank (lemma 9 of IV, p. 297) and the one relative to \(\{\lambda\}\) is assumed to have finite rank, it follows from prop. 8 of IV, p. 278 that there exists a sequence \((k_n)_{n \in \mathbf{N}}\) strictly increasing in \(\mathbf{N}\) such that the spectral projector \(p_{n}\) of u relative to \(J_{k_{n}}\) is nonzero. Let \(x_{n}\) be a vector of norm 1 in the image of \(p_{n}\) . The sequence \((x_{n})\) is orthonormal, since the image of \(p_{n}\) is orthogonal to that of \(p_{m}\) for all \(n \neq m\) in N.

Let \(n \in N\) . Denote by \(\mu_{n}\) the spectral measure of \(x_{n}\) relative to u; its support is contained in \(J_{k_{n}}\) (prop. 9 of IV, p. 278). It follows

\[
\| u (x _ {n}) - \lambda x _ {n} \| ^ {2} = \int_ {\mathbf {C}} | t - \lambda | ^ {2} d \mu_ {n} (t) \leqslant \frac {1}{k _ {n} ^ {2}},
\]

therefore the sequence \((x_{n})_{n\in\mathbf{N}}\) has the required property.

Suppose conversely that there exists an orthonormal sequence \((x_{n})_{n\in \mathbf{N}}\) in E such that \(u(x_{n}) - \lambda x_{n}\to 0\) . Denote by \(\mu_{n}\) the spectral measure of \(x_{n}\) relative to \(u\) .

Let \(\varepsilon > 0\) . Denote by \(p_{\varepsilon}\) be the spectral projector of \(u\) relative to the open interval \(\mathrm{I}_{\varepsilon} = ]\lambda - \varepsilon, \lambda + \varepsilon[\) . For every \(n \in \mathbf{N}\) , it follows that

\[
\begin{array}{r l} & 1 = \| x _ {n} \| ^ {2} = \mu_ {n} (\mathrm{I} _ {\varepsilon}) + \mu_ {n} (\mathbf {R} - \mathrm{I} _ {\varepsilon}) \\ & \qquad \leqslant \mu_ {n} (\mathrm{I} _ {\varepsilon}) + \frac {1}{\varepsilon^ {2}} \int_ {\mathbf {R} - \mathrm{I} _ {\varepsilon}} | t - \lambda | ^ {2} d \mu_ {n} (t) \\ & \qquad = \| p _ {\varepsilon} (x _ {n}) \| ^ {2} + \frac {1}{\varepsilon^ {2}} \| u (x _ {n}) - \lambda x _ {n} \| ^ {2}. \end{array}
\]

The hypothesis on the sequence \((x_{n})\) therefore implies that the norm of \(p_{\varepsilon}(x_{n})\) cannot tend to 0. Since the orthonormal sequence \((x_{n})\) converges weakly to 0 in E (Lemma 13), the projector \(p_{\varepsilon}\) cannot be compact (cor. of prop. 6 of III, p. 6) and is consequently of infinite rank. Since this holds for every \(\varepsilon > 0\) , lemma 9 of IV, p. 297 allows us to conclude that \(\lambda \in \mathrm{Sp}_{\mathrm{e}}(u)\) .

The following theorem is the analogue for self-adjoint partial operators of Theorem 3 in III, p. 93.

THEOREM 4. --- Let u be a self-adjoint partial operator on E. The essential spectrum of u is the intersection of the sets Sp(u + v) where v ranges over the set of compact Hermitian endomorphisms of E.

The partial operator \(u + v\) is self-adjoint for every Hermitian endomorphism v of E since \((u + v)^{*} = u^{*} + v^{*}\) (cf. IV, p. 236).

Let us prove that if \(v\) is compact, then \(\mathrm{Sp_e}(u) \subset \mathrm{Sp_e}(u + v)\) . Let \(\lambda \in \mathrm{Sp_e}(u)\) and let \((x_n)_{n \in \mathbf{N}}\) be an orthonormal sequence in E such that \(u(x_n) - \lambda x_n\) converges to 0 (prop. 18). The sequence \((x_n)\) converges weakly to 0 in E (lemma 13), and since the endomorphism \(v\) is compact, the sequence \((v(x_n))_{n \in \mathbf{N}}\) converges to 0 in E (cor. of prop. 6 of III, p. 6). Consequently, the sequence \((u+v)(x_{n})-\lambda x_{n}\) converges to 0 in E, and Prop. 18 implies that \(\lambda\in\mathrm{Sp}_{\mathrm{e}}(u+v)\) .

The essential spectrum of \(u\) is therefore contained in the intersection of the sets \(\mathrm{Sp}(u + v)\) for Hermitian \(v \in \mathcal{L}^{\mathrm{c}}(\mathrm{E})\).

Conversely, let \(\lambda \in \mathrm{Sp}_{\mathrm{s}}(u)\) . Denote by \(\mathrm{E}_{\lambda}\) the eigenspace of \(u\) relative to \(\lambda\) , and \(p_{\lambda}\) the orthogonal projector with image \(\mathrm{E}_{\lambda}\) ; it is the spectral projector of \(u\) relative to \(\{\lambda\}\) (corollary of Prop. 9 of IV, p. 278). By definition of the sensitive spectrum, the projector \(p_{\lambda}\) is of finite rank, hence compact. Set \(w = u + p_{\lambda}\) ; it is a self-adjoint partial operator. To conclude the proof, let us verify that \(\lambda\) belongs to the resolvent set of \(w\) .

One has \(\mathrm{E}_{\lambda} \subset \operatorname{dom}(u)\) and the spaces \(\mathrm{E}_{\lambda}\) and \(\operatorname{dom}(u) \cap \mathrm{E}_{\lambda}^{\circ}\) are stable under \(u\) (prop. 9 of IV, p. 278).

Let \(x \in \text{dom}(u)\) . Let us write \(x = p_{\lambda}(x) + y\) where \(y \in \text{dom}(u) \cap \text{E}_{\lambda}^{\circ}\) . From what precedes, we have

\[
\left\| w (x) - \lambda x \right\| ^ {2} = \left\| p _ {\lambda} (x) \right\| ^ {2} + \left\| w (y) - \lambda y \right\| ^ {2}.\tag{21}
\]

Let V be an open neighborhood of \(\lambda\) not meeting the spectrum of u, and let c > 0 be such that the disk with center \(\lambda\) and of radius c is contained in V. Let \(\mu_{y}\) be the spectral measure of y relative to u. Since y is orthogonal to \(E_{\lambda}\) , the support of \(\mu_{y}\) does not meet V (loc. cit.). We then have

\[
\| w (y) - \lambda y \| ^ {2} = \| u (y) - \lambda y \| ^ {2} = \int_ {\mathbf {C}} | t - \lambda | ^ {2} d \mu_ {y} (t) \geqslant c ^ {2} \| y \| ^ {2}.\tag{22}
\]

We conclude from (21) and (22) that \(\|w(x)-\lambda x\|^{2}\geqslant\inf(c^{2},1)\|x\|^{2}\), since w is self-adjoint and \(\lambda\in R\). The conclusion then follows from Lemma 5 of IV, p. 248.

COROLLARY. --- Let \(u\) be a self-adjoint partial operator on E and \(v\) a compact Hermitian endomorphism of E. We have \(\mathrm{Sp}_{\mathrm{e}}(u + v) = \mathrm{Sp}_{\mathrm{e}}(u)\) .

This follows immediately from the theorem.

\subsection{Perturbation}

In this number, E is a complex Hilbert space.

If \(u\) is a self-adjoint partial operator on E and \(v \in \mathcal{L}(\mathrm{E})\) is a Hermitian endomorphism, then \(u + v\) is self-adjoint (cf. IV, p. 236). This is not the case in general when \(v\) is a symmetric partial operator (exercise 9 of IV, p. 347). We shall nevertheless obtain positive results of this type in this number.

DEFINITION 11. --- Let u be a partial operator on E. A partial operator v on E is said to be bounded relative to u if dom(u) ⊂ dom(v) and if v defines, by restriction to the subspace, a continuous linear map from \(E_{u}\) into E.

Let u be a partial operator on E. Let v be a partial operator bounded relative to u. There exists a real number m such that

\[
\| v (x) \| \leqslant m (\| x \| + \| u (x) \|)
\]

for all \(x \in \text{dom}(u)\) . The relative norm of v with respect to u, denoted by \(\|v\|_{u}\), is the lower bound of the real numbers \(a \geqslant 0\) such that there exists a real number b such that

\[
\| v (x) \| \leqslant a \| u (x) \| + b \| x \|
\]

for all \(x \in \text{dom}(u)\) . The relative norm of v is therefore less than the norm of the restriction of v to the space \(E_{u}\) .

Note. --- Let u be a partial operator on E. Every endomorphism \(v \in \mathcal{L}(\mathrm{E})\) is bounded relative to u and its relative norm is zero since \(\|v(x)\| \leqslant \|v\|\|x\|\) for all \(x \in \operatorname{dom}(u)\) , which allows us to take a = 0 in the inequality above.

THEOREM 5 (Kato–Rellich). — Let u be a self-adjoint partial operator on E and v a symmetric partial operator bounded relative to u. If the relative norm \(\|v\|_{u}\) is \(< 1\), then the partial operator \(u + v\) whose domain is dom(u) is self-adjoint.

Since \(\| v \|_u < 1\) there exist, by definition, positive real numbers \(a\) and \(b\) such that \(a < 1\) and \(\| v(x) \| \leqslant a \| u(x) \| + b \| x \|\) for all \(x \in \text{dom}(u)\) .

Let \(t \in R^{*}\) . Set \(w_{t} = v \circ \mathrm{R}(u, it)\) . We have \(\operatorname{dom}(w_{t}) = E\) since the image of \(\mathrm{R}(u, it)\) is contained in the domain of u, which is contained in the domain of v by hypothesis. Let \(x \in \operatorname{dom}(u)\) . Since u is self-adjoint, we have

\[
\begin{array}{c} \| (i t - u) x \| ^ {2} = \| i t x \| ^ {2} + \| u (x) \| ^ {2} - i t \langle x | u (x) \rangle + i t \langle u (x) | x \rangle \\ = | t | ^ {2} \| x \| ^ {2} + \| u (x) \| ^ {2}, \end{array}
\]

whence the inequalities \(\|u(x)\|\leqslant\|(it-u)x\|\) and \(\|x\|\leqslant|t|^{-1}\|(it-u)x\|\) . But then it follows that

\[
\| v (x) \| \leqslant (a + b | t | ^ {- 1}) \| (i t - u) x \|.
\]

In particular, set \(x = \mathrm{R}(u, it)y\) for \(y \in \mathrm{E}\) ; we obtain

\[
\| w _ {t} (y) \| \leqslant (a + b | t | ^ {- 1}) \| y \|.
\]

We deduce that \(w_{t}\in\mathcal{L}(\mathrm{E})\) and that \(\|w_{t}\|\leqslant a+|t|^{-1}b\) .

Since \(a < 1\) , there exists \(t \in \mathbf{R}_{+}^{*}\) such that \(a + b|t|^{-1} < 1\) . We then have \(\|w_{t}\| < 1\) and \(\|w_{-t}\| < 1\) ; the endomorphisms \(1_{\mathrm{E}} - w_{t}\) and \(1_{\mathrm{E}} - w_{-t}\) of E are therefore invertible (prop. 2 of I, p. 22). Since we have the formula \((1_{\mathrm{E}} - w_{t}) \circ (it - u) = it - (u + v)\) , this implies that \(u + v - it\) is surjective; likewise, the partial operator \(u + v + it\) is surjective. It follows that \(u + v\) is self-adjoint (prop. 11 of IV, p. 240).

\subsection{Operators with compact resolvent}

In this number, E denotes a complex Hilbert space of infinite dimension.

PROPOSITION 19. --- Let u be a self-adjoint partial operator on E. The following conditions are equivalent:

\begin{enumerate}
\def\labelenumi{(\roman{enumi})}
\item
  There exists an orthonormal basis \(\mathrm{B} = (e_{j})_{j \in \mathrm{J}}\) of E such that u is diagonal in the basis B and the absolute value of the family of eigenvalues of u tends to infinity along the filter of complements of finite subsets of J;
\item
  For every \(\lambda\) belonging to the resolvent set of \(u\) , the resolvent \(\mathrm{R}(u,\lambda)\) is compact;
\item
  There exists a complex number \(\lambda\) belonging to the resolvent set of \(u\) such that the resolvent \(\mathrm{R}(u,\lambda)\) is compact;
\item
  The spectrum of u coincides with the sensitive spectrum of u.
\end{enumerate}

Suppose that (i) is satisfied and let \((\lambda_{j})_{j\in\mathrm{J}}\) be the family of eigenvalues of u. Let \(\mu\) be the counting measure on J. The spectrum of u is the image \(\mu\) -essential of the family \((\lambda_{j})\) (prop. 22 of IV, p. 252); it is the set of values of this family. For every \(\lambda\notin\mathrm{Sp}(u)\) , the resolvent \(\mathrm{R}(u,\lambda)\) is diagonal in the basis B and the family of its eigenvalues is \(((\lambda-\lambda_{j})^{-1})_{j\in\mathrm{J}}\) (loc. cit.). Since this family converges to 0, the endomorphism \(\mathrm{R}(u,\lambda)\) is compact (prop. 2, (iii) of IV, p. 148). This proves that (i) implies (ii).

Since the resolvent set of \(u\) is nonempty (cf. prop. 17 of IV, p. 248), properties (ii) and (iii) are equivalent by formula (8) of IV, p. 245 and proposition 3 of III, p. 5.

Since \(\mathrm{R}(u,\lambda)\) is normal for \(\lambda\notin\mathrm{Sp}(u)\) (prop. 16 of IV, p. 247), condition (iii) implies (iv) by prop. 15 of IV, p. 247 and prop. 5 of III, p. 90.

Finally, let us prove that (iv) implies (i). Let \(\mathcal{O}\) be the set of orthonormal subsets of E consisting of eigenvectors for \(u\) . The set \(\mathcal{O}\) , ordered by inclusion, is of finite character (E, III, p. 34, def. 2) since O belongs to \(\mathcal{O}\) if and only if the sets formed of at most two elements of O belong to \(\mathcal{O}\) . By E, III, p. 35, th. 1, there exists a maximal element O of \(\mathcal{O}\) . Let F be the closed subspace of E generated by O. For \(e \in O\) , there exists a unique \(\lambda(e) \in \mathbf{R}\) such that \(u(e) = \lambda(e)e\) (prop. 17 of IV, p. 248).

The set of values of the map \(\lambda\) from O to \(\mathbf{R}\) coincides with the spectrum of \(u\) . Indeed, on the one hand this set is contained in the spectrum of \(u\) and on the other hand, if there exists \(\lambda_0 \in \mathrm{Sp}(u)\) which is not a value of \(\lambda\) , then \(\lambda_0\) is an eigenvalue of \(u\) by hypothesis. There exists a vector \(e \in \mathrm{dom}(u)\) of norm 1 such that \(u(e) = \lambda_0 e\) and \(e \notin \mathrm{O}\) ; the subset \(\mathrm{O} \cup \{e\}\) is orthonormal (since eigenspaces of \(u\) corresponding to distinct eigenvalues are orthogonal by assertion \(b\) ) of the cor. of prop. 17 of IV, p. 248); it belongs to \(\mathcal{O}\) , contradicting the maximality of O.

Suppose that \(F \neq E\) . Then \(F^{\circ}\) is nonzero. The endomorphism \(\mathrm{R}(u,i)\) of E is normal (prop. 16 of IV, p. 247). It leaves the subspace \(F^{\circ}\) of E stable (lemma 4 of I, p. 135). Let v be the endomorphism of \(F^{\circ}\) deduced from \(\mathrm{R}(u,i)\) by passing to subspaces. It is a continuous and normal endomorphism of \(F^{\circ}\) (loc. cit.) whose spectrum is contained in that of \(\mathrm{R}(u,i)\) . Since \(F^{\circ}\) is nonzero, the spectrum of v is nonempty (cor. 1 of I, p. 26). Moreover, the spectrum of v is not reduced to 0, since the normal endomorphism v is nonzero (example 1 of I, p. 110). Let \(s \in \mathrm{Sp}(v) - \{0\}\) . Since s belongs to the spectrum of \(\mathrm{R}(u,i)\) , there exists \(e \in O\) such that \(s = (i - \lambda(e))^{-1}\) , and s is an eigenvalue of \(\mathrm{R}(u,i)\) (prop. 15, a) of IV, p. 247). Since s is nonzero, it is an isolated point of \(\mathrm{Sp}(\mathrm{R}(u,i))\) by hypothesis, hence also of \(\mathrm{Sp}(v)\) . Thus, s is an eigenvalue of v (prop. 5, c) of I, p. 134). Let \(e \in F^{\circ}\) be an eigenvector of norm 1 of v; it is also an eigenvector of u (prop. 15, b) of IV, p. 247), and the set \(O \cup \{e\}\) contradicts the fact that O is maximal in O. Hence F = E.

The family of elements of O is therefore an orthonormal basis of E consisting of eigenvectors of u. The spectrum of u is closed in R, and the sensitive spectrum is discrete; since these sets coincide, the set of elements \(\lambda \in \mathrm{Sp}(u)\) such that \(|\lambda| \leqslant R\) is compact, hence finite, for every R > 0. Thus the family of absolute values of the eigenvalues of u tends to infinity along the filter of complements of finite subsets of O. Hence (iv) implies (i).

DEFINITION 12. --- Let u be a self-adjoint partial operator on E. One says that u has compact resolvent if the equivalent conditions of Prop. 19 are satisfied.

PROPOSITION 20. --- Let u be a self-adjoint partial operator on E. Then u has compact resolvent if and only if the canonical injection j of the Hilbert space \(E_{u}\) into E is compact.

Suppose that u has compact resolvent. There exists \(\lambda \notin \operatorname{Sp}(u)\) such that \(\mathrm{R}(u, \lambda)\) is compact (prop. 19, (iii)). Let B be the unit ball of the Hilbert space \(E_{u}\) . Since u is a continuous linear map from \(E_{u}\) into E, the subset \(\mathrm{B}' = (\lambda 1_{\mathrm{E}} - u)(\mathrm{B})\) of E is bounded. Since \(\mathrm{R}(u, \lambda)\) is compact, the subset \(\mathrm{B} = \mathrm{R}(u, \lambda)(\mathrm{B}')\) is relatively compact in E (Remark 1 of III, p. 2). This proves that j is compact.

Conversely, suppose that the linear map \(j\) is compact. The complex number \(i\) belongs to the resolvent set of \(u\) (prop. 17 of IV, p. 248). We have \(u \circ \mathrm{R}(u, i) = -1_{\mathrm{E}} + i\mathrm{R}(u, i)\) . Let B be the unit ball of E. For every \(x \in \mathrm{B}\) , we have

\[
\| u \circ \mathrm{R} (u, i) (x) \| = \| - x + i \mathrm{R} (u, i) (x) \| \leqslant 1 + \| \mathrm{R} (u, i) \|,
\]

and consequently

\[
\begin{array}{r l} & {\| \mathrm{R} (u, i) x \| _ {u} ^ {2} = \| \mathrm{R} (u, i) x \| ^ {2} + \| u \circ \mathrm{R} (u, i) x \| ^ {2}} \\ & {\qquad \leqslant \| \mathrm{R} (u, i) \| ^ {2} + (1 + \| \mathrm{R} (u, i) \|) ^ {2}.} \end{array}
\]

The image C of B under \(R(u,i)\) is therefore bounded in \(E_{u}\) . Since j is compact by hypothesis, the set \(C = j(C)\) is relatively compact in E (III, loc. cit.). Consequently, the resolvent \(R(u,i)\) is compact and u has compact resolvent (prop. 19, (iii)).

COROLLARY. --- Let q be a closed positive partial form on E and u the positive self-adjoint operator representing q. The following conditions are equivalent:

\begin{enumerate}
\def\labelenumi{(\roman{enumi})}
\item
  The partial operator u has compact resolvent;
\item
  The positive endomorphism \(\sqrt{(1_{\mathrm{E}} + u)^{-1}} = (1_{\mathrm{E}} + u)^{-1 / 2}\) of E is compact;
\item
  The canonical injection j of the Hilbertian space \(E_{q}\) (def. 9 of IV, p. 292) in E is compact.
\end{enumerate}

Since u is positive, the real number -1 belongs to the resolvent set of u (prop. 17 of IV, p. 248), hence the positive endomorphism \(v = \sqrt{(1_{\mathrm{E}} + u)^{-1}}\) is defined, and we have \(v = (1_{\mathrm{E}} + u)^{-1/2}\) by the functional calculus.

The endomorphism \(v\) is compact if and only if \(v^2 = (1_{\mathrm{E}} + u)^{-1}\) is compact (prop. 6 of III, p. 91), that is, if and only if \(u\) has compact resolvent (prop. 19, (iii)). This proves that conditions (i) and (ii) are equivalent.

Suppose that the endomorphism \(v\) is compact. Let B be the unit ball of the Hilbert space \(\mathrm{E}_q\) . Since \((1_{\mathrm{E}} + u)^{1/2}\) is a continuous linear map from \(\mathrm{E}_q\) into E (cor. of th. 3 of IV, p. 294), the subset \(\mathrm{B}' = (1_{\mathrm{E}} + u)^{1/2}(\mathrm{B})\) of E is bounded, hence the subset \(j(\mathrm{B}) = \mathrm{B} = v(\mathrm{B}')\) is relatively compact in E (Remark 1 of III, p. 2). This proves that \(j\) is compact. Thus (ii) implies (iii).

The linear map \(\widetilde{v}\colon x\mapsto(1_{\mathrm{E}}+u)^{-1/2}(x)\) from E into \(\mathrm{E}_{q}\) is well defined and isometric (corollary of Th. 3 of IV, p. 294). Since \(v=j\circ\widetilde{v}\) , condition (iii) implies (ii) (Prop. 3 of III, p. 5).

Example. --- Let \(n \in N\) . Let U be an open subset of \(R^{n}\) . We equip U with Lebesgue measure, denoted \(\mu\) . Let \(\Delta\) be the scalar differential operator \(-\sum_{i=1}^{n}\partial_{i}^{2}\) on U. The partial operator \(\Delta_{-}\) with domain \(\mathcal{D}(U)\) defined by \(\varphi \mapsto \Delta(\varphi)\) is closable (prop. 13 of IV, p. 242) and symmetric (IV, p. 243). It is positive, since for every \(\varphi \in \mathcal{D}(U)\) , we have

\[
\begin{array}{r l} & {\langle \varphi | \Delta_ {-} (\varphi) \rangle = \int_ {\mathrm{U}} \overline {{\varphi}} \Delta (\varphi) d \mu = - \sum_ {i = 1} ^ {n} \int_ {\mathrm{U}} \overline {{\varphi}} \partial_ {i} ^ {2} \varphi d \mu} \\ & {\qquad = \sum_ {i = 1} ^ {n} \int_ {\mathrm{U}} \overline {{\partial_ {i} \varphi}} \partial_ {i} \varphi d \mu = \int_ {\mathrm{U}} \sum_ {i = 1} ^ {n} | \partial_ {i} \varphi | ^ {2} d \mu \geqslant 0.} \end{array}
\]

We denote by \(\Delta_{D}\) the Friedrichs extension of the positive symmetric partial operator \(\Delta_{-}\) (IV, p. 295, example); it is a Laplacian on U, called the Dirichlet Laplacian on U.

Let q be the positive partial form associated with \(\Delta_{D}\) . The domain of q is the Hilbertian space completion of \(\mathcal{D}(U)\) for the positive-definite Hermitian form

defined by

\[
(\varphi_ {1}, \varphi_ {2}) \mapsto \int_ {\mathrm{U}} \overline {{\varphi}} _ {1} \varphi_ {2} + \sum_ {i = 1} ^ {n} \int_ {\mathrm{U}} \overline {{\partial_ {i} \varphi_ {1}}} \partial_ {i} \varphi_ {2}
\]

for all \((\varphi_{1},\varphi_{2})\in\mathcal{D}(\mathrm{U})\times\mathcal{D}(\mathrm{U})\) . In other words, the domain of \(q\) is the Sobolev space \(\mathrm{H}_{0}^{1}(\mathrm{U})\) (No. 14 of IV, p. 221).

*Suppose that U is bounded. The canonical injection of \(\mathrm{H}_0^1 (\mathrm{U})\) into \(\mathrm{L}^2 (\mathrm{U})\) is compact; the Dirichlet Laplacian on U is therefore an operator with compact resolvent (corollary of Prop. 20). Since the Hilbert space \(\mathrm{H}_0^1 (\mathrm{U})\) is of countable type (prop. 20 of IV, p. 222), and since the image of \(\Delta_{\mathrm{D}}\) is of infinite dimension, there exists an increasing sequence \((\lambda_n)_{n\geqslant 0}\) of real numbers tending to \(+\infty\) and an orthonormal basis \((f_n)_{n\in \mathbf{N}}\) of \(\mathrm{L}^2 (\mathrm{U})\) whose elements belong to the domain of \(\Delta_{\mathrm{D}}\) , such that \(\Delta_{\mathrm{D}}(f_n) = \lambda_n f_n\) for all \(n\in \mathbf{N}\) . One can prove (“Weyl’s law”) that when T tends to \(+\infty\) , we have

\[
\sum_{\substack{n\in \mathbf{N}\\ \lambda_{n}\leqslant \mathrm{T}}}1\sim \frac{c_{n}}{(2\pi)^{n}} m\mathrm{T}^{n / 2}
\]

where \(c_{n} = \pi^{n/2}/\Gamma(1 + n/2)\) is the volume of the unit ball in \(R^{n}\) (INT, V, p. 101, § 8, no. 7) and \(m > 0\) is the Lebesgue measure of U.*

\specialchapter{Exercises}

\exercisesection{}

\begin{enumerate}
\def\labelenumi{\arabic{enumi})}
\tightlist
\item
  Let \(E = \ell_{\mathbf{C}}^{2}(\mathbf{N})\) and let \(u \in \mathcal{L}(\mathrm{E})\) be the endomorphism such that \(u(x_{0}, \ldots, x_{n}, \ldots) = (0, x_{0}, \ldots, x_{n}, \ldots)\) for all \((x_{n})_{n \in \mathbf{N}}\) in E. Let \((e_{n})\) be the canonical orthonormal basis of E.
\end{enumerate}

\begin{enumerate}
\def\labelenumi{\alph{enumi})}
\item
  For every \(v \in \mathcal{L}(\mathrm{E})\) such that \(\|v\| < 1/2\) , the vector \(e_{0}\) is not adherent to the image of \(u + v\) .
\item
  The open ball in \(\mathcal{L}(\mathrm{E})\) of radius 1/2 centered at u contains no diagonalizable endomorphism of E.
\end{enumerate}

\begin{enumerate}
\def\labelenumi{\arabic{enumi})}
\setcounter{enumi}{1}
\item
  Let E be a complex Hilbert space. Prove directly that an endomorphism \(u\) of E is a Fredholm endomorphism if and only if \(\pi(u)\) is invertible in \(\mathcal{C}alk(E)\) (Atkinson's theorem). (Use the spectral properties of compact operators, cf. § 1 of IV, p. 146.)
\item
  Let E be a complex Hilbert space and let \(\mathrm{B}=(e_{i})_{i\in\mathrm{I}}\) be an orthonormal basis of E. For every \(i\in I\) , we denote by \(p_{i}\) the orthogonal projector with image \(Ce_{i}\) .
\end{enumerate}

\begin{enumerate}
\def\labelenumi{\alph{enumi})}
\item
  Let K be a precompact subset of E. For every \(\varepsilon > 0\) there exists a finite subset J of I such that, for every \(x \in \mathbf{K}\) , \(\|x - \sum_{i \in \mathbf{J}} p_j(x)\| \leqslant \varepsilon\) .
\item
  Let \(u \in \mathcal{D}_{\mathrm{B}}(\mathrm{E})\) and let us denote by \(\lambda = (\lambda_{i})_{i \in \mathrm{I}}\) its family of eigenvalues. The family \((\lambda_{i} p_{i})\) is summable in \(\mathcal{L}(\mathrm{E})\) endowed with the topology of precompact convergence and whose sum is equal to u (cf. prop. 1 of IV, p. 147)
\end{enumerate}

\begin{enumerate}
\def\labelenumi{\arabic{enumi})}
\setcounter{enumi}{3}
\item
  Let E be a complex Hilbert space of infinite dimension. There exists a compact endomorphism of E that is not of finite trace.
\item
  \begin{enumerate}
  \def\labelenumii{\alph{enumii})}
  \tightlist
  \item
  Let \(E = L^{2}([0,1], \mu)\) where \(\mu\) is the Lebesgue measure. The image A of the algebra \(L^{\infty}([0,1], \mu)\) by the map that associates to f the endomorphism of multiplication by f is a maximal commutative subalgebra of \(\mathcal{L}(E)\) which is not of the form \(\mathcal{D}_{\mathrm{B}}(\mathrm{E})\) for an orthonormal basis B of E. (Verify that the algebra A contains no minimal projector.)
  \end{enumerate}
\end{enumerate}

\begin{enumerate}
\def\labelenumi{\alph{enumi})}
\setcounter{enumi}{1}
\tightlist
\item
  Let E be an infinite-dimensional Hilbert space. There exists a maximal commutative subalgebra of \(\mathcal{L}(\mathrm{E})\) which is not of the form \(\mathcal{D}_{\mathrm{B}}(\mathrm{E})\) for B an orthonormal basis of E.
\end{enumerate}

\begin{enumerate}
\def\labelenumi{\arabic{enumi})}
\setcounter{enumi}{5}
\tightlist
\item
  Let E be a complex Hilbert space.
\end{enumerate}

\begin{enumerate}
\def\labelenumi{\alph{enumi})}
\item
  Let u be a compact normal endomorphism of E. Every closed subspace of E invariant under u is invariant under \(u^{*}\) .
\item
  If E is infinite-dimensional, there exist a normal endomorphism u of E and a closed subspace \(F \subset E\) such that:
\end{enumerate}

\begin{enumerate}
\def\labelenumi{(\roman{enumi})}
\item
  One has \(u(\mathrm{F}) \subset \mathrm{F}\) ;
\item
  The space F is not invariant under \(u^{*}\) ;
\item
  The space F° is not invariant under u ;
\item
  The endomorphism of F induced by \(u\) by passing to subspaces is not normal.
\end{enumerate}

\begin{enumerate}
\def\labelenumi{\arabic{enumi})}
\setcounter{enumi}{6}
\tightlist
\item
  The goal of this exercise is to give a more direct proof of Theorem 1 of IV, p. 149. Let E be an infinite-dimensional complex Hilbert space and let u be a compact and normal endomorphism of E.
\end{enumerate}

\begin{enumerate}
\def\labelenumi{\alph{enumi})}
\item
  Show that, for every nonzero complex number \(\lambda\), the eigenspace \(\operatorname{Ker}(u - \lambda 1_{\mathrm{E}})\) is of finite dimension.
\item
  Suppose that \(u\) is Hermitian. Show that there exists \(x_0 \neq 0\) in E such that
\end{enumerate}

\[
\frac {\| u (x _ {0}) \| ^ {2}}{\| x _ {0} \| ^ {2}} = \sup _ {x \in \mathrm{E} - \{0 \}} \frac {\| u (x) \| ^ {2}}{\| x \| ^ {2}}.
\]

\begin{enumerate}
\def\labelenumi{\alph{enumi})}
\setcounter{enumi}{2}
\item
  Show that \(x_{0}\) is an eigenvector of \(u\) for an eigenvalue \(\lambda\) such that \(|\lambda| = \|u\|\) .
\item
  Show that there exists a countable set I, an orthonormal family \((e_{i})_{i\in\mathrm{I}}\) in E and a family \((\lambda_{i})_{i\in\mathrm{I}}\) of real numbers such that
\end{enumerate}

\[
u (x) = \sum_ {i \in I} \lambda_ {i} \langle e _ {i} | x \rangle e _ {i}
\]

for all \(x \in E\) .

\begin{enumerate}
\def\labelenumi{\alph{enumi})}
\setcounter{enumi}{4}
\tightlist
\item
  Extend the result to the case where u is compact and normal. (Write \(u = u_{1} + iu_{2}\) where \(u_{1}\) and \(u_{2}\) are compact, Hermitian, and mutually commuting.)
\end{enumerate}

\begin{enumerate}
\def\labelenumi{\arabic{enumi})}
\setcounter{enumi}{7}
\tightlist
\item
  We endow the group R/Z with its normalized Haar measure \(\mu\) .
\end{enumerate}

\begin{enumerate}
\def\labelenumi{\alph{enumi})}
\tightlist
\item
  Let \(h \in \mathcal{L}^{2}(\mathbf{R}/\mathbf{Z})\) . For \(f \in \mathcal{L}^{2}(\mathbf{R}/\mathbf{Z})\) we define a function \(\mathrm{C}_{h}(f)\) on R/Z by the formula
\end{enumerate}

\[
\mathrm{C} _ {h} (f) (x) = \int_ {\mathbf {R} / \mathbf {Z}} f (t) h (x - t) d \mu (t).
\]

The map \(f \mapsto \mathrm{C}_h(f)\) is a continuous endomorphism of \(\mathcal{L}^2(\mathbf{R}/\mathbf{Z})\) , and by passing to quotients defines a compact normal endomorphism of \(\mathrm{L}^2(\mathbf{R}/\mathbf{Z})\) .

\begin{enumerate}
\def\labelenumi{\alph{enumi})}
\setcounter{enumi}{1}
\item
  Determine the spectrum of \(C_{h}\) and the spectral multiplicity of the elements of this spectrum. (Expand h in a Fourier series.)
\item
  There exists \(h \in \mathcal{C}(\mathbf{R}/\mathbf{Z})\) such that \(C_{h}\) is not of finite trace.
\end{enumerate}

\begin{enumerate}
\def\labelenumi{\arabic{enumi})}
\setcounter{enumi}{8}
\tightlist
\item
  We denote by E the complex Hilbert space \(L^{2}([0,1])\) of functions that are square-integrable on \([0,1]\) with respect to Lebesgue measure \(\mu\) .
\end{enumerate}

Let \(u_{0} \in \mathcal{L}(E)\) be the Volterra endomorphism on E, such that

\[
u _ {0} (f) (x) = \int_ {[ 0, x ]} f (t) d \mu (t)
\]

for all \(f \in \mathrm{L}^2([0,1])\) and almost all \(x \in [0,1]\) .

\begin{enumerate}
\def\labelenumi{\alph{enumi})}
\item
  The endomorphism \(u_{0}\) is compact and \(\operatorname{Sp}(u_{0})=\{0\}\) .
\item
  The endomorphism \(u_0^* u_0\) has finite trace; compute its trace and deduce the value of the series
\end{enumerate}

\[
\sum_ {n = 1} ^ {+ \infty} \frac {1}{n ^ {2}}.
\]

\begin{enumerate}
\def\labelenumi{\alph{enumi})}
\setcounter{enumi}{2}
\tightlist
\item
  Let \(v_{0}\) be the rank-1 endomorphism of E defined by
\end{enumerate}

\[
v _ {0} (f) = \left(\int_ {0} ^ {1} f (t) d \mu (t)\right) \cdot 1,
\]

where 1 denotes the class of the constant function 1. Let \(u = u_0 - v_0\) . Show that \(u\) is a Hilbert--Schmidt operator, and compute the kernel N such that

\[
u (f) (x) = \int_ {0} ^ {1} \mathrm{N} (x, y) f (y) d \mu (y)
\]

for \(f \in \mathrm{E}\) and almost all \(x \in [0,1]\) .

\begin{enumerate}
\def\labelenumi{\alph{enumi})}
\setcounter{enumi}{3}
\item
  Determine the eigenvalues of u and their spectral multiplicity. (Compute \(u(e_{n})\) , where \(e_{n}(t) = e^{2i\pi nt}\) for all \(n \in Z\) .)
\item
  Compute the trace of \(u^{*}u\) and deduce anew the value of the series whose general term is \(1/n^{2}\) for \(n \geqslant 1\) .
\item
  Generalize these arguments to show that, if \(k \geqslant 1\) is a strictly positive integer, then
\end{enumerate}

\[
\sum_ {n = 1} ^ {+ \infty} \frac {1}{n ^ {2 k}} \in \pi^ {2 k} {\bf Q} ^ {*}.
\]

\begin{enumerate}
\def\labelenumi{\arabic{enumi})}
\setcounter{enumi}{9}
\tightlist
\item
  Let E be a complex Hilbert space. Let A be a star-subalgebra of \(\mathcal{L}(\mathrm{E})\) . Suppose that the representation of A in E is irreducible and that A contains a nonzero compact endomorphism.
\end{enumerate}

\begin{enumerate}
\def\labelenumi{\alph{enumi})}
\item
  The star-subalgebra A contains a nonzero finite-rank projector. (Show that A contains a nonzero compact Hermitian endomorphism and apply the functional calculus.)
\item
  The star-subalgebra A contains a rank-1 projector. (If \(p \in \mathrm{A}\) is a nonzero projector of minimal rank, show that the algebra \(p\mathrm{A}p\) is contained in \(\mathbf{C}p\) .)
\item
  It follows that A contains \(\mathcal{L}^{\mathrm{c}}(\mathrm{E})\) .
\item
  Conversely, every star-subalgebra of \(\mathcal{L}(\mathrm{E})\) containing \(\mathcal{L}^{\mathrm{c}}(\mathrm{E})\) is irreducible.
\end{enumerate}

\begin{enumerate}
\def\labelenumi{\arabic{enumi})}
\setcounter{enumi}{10}
\tightlist
\item
  Let E be a complex Hilbert space. Let \(A \subset \mathcal{L}(\mathrm{E})\) be a set of compact endomorphisms such that \(u \circ v = v \circ u\) for every pair \((u, v)\) of elements of A and \(u^{*} \in A\) for all \(u \in A\) .
\end{enumerate}

Show that there exists a countable orthonormal family \((e_i)_{i\in \mathrm{I}}\) and a family \((\Lambda_i)_{i\in \mathrm{I}}\) of mappings \(\mathrm{A}\to \mathbf{C}\) such that

\[
u (x) = \sum_ {i \in \mathrm{I}} \Lambda_ {i} (u) \langle e _ {i} | x \rangle e _ {i}
\]

for every element u of A and every x in E. If A is a unital subalgebra of \(\mathcal{L}(\mathrm{E})\) , then \(\Lambda_{i}\) is a character of A for every \(i \in I\) .

\begin{enumerate}
\def\labelenumi{\arabic{enumi})}
\setcounter{enumi}{11}
\tightlist
\item
  If E is a complex Hilbert space and u is a compact normal endomorphism of E, show that u has a cyclic vector (cf. def. 3 of IV, p. 191) if and only if the following conditions are satisfied:
\end{enumerate}

\begin{enumerate}
\def\labelenumi{(\roman{enumi})}
\item
  The kernel of u is reduced to 0;
\item
  The spectral multiplicity of every nonzero eigenvalue of u is equal to 1.
\end{enumerate}

\begin{enumerate}
\def\labelenumi{\arabic{enumi})}
\setcounter{enumi}{12}
\tightlist
\item
  Let E be a complex Hilbert space of countable type, let \((e_{n})_{n\in\mathbf{N}}\) be an orthonormal basis of E, and let u be an endomorphism of E. Show that u is compact if and only if one has
\end{enumerate}

\[
\lim _ {n \to + \infty} \sup _ {x \in \{e _ {0}, \dots , e _ {n - 1} \} ^ {\circ} - \{0 \}} \frac {\| u (x) \| ^ {2}}{\| x \| ^ {2}} = 0.
\]

\begin{enumerate}
\def\labelenumi{\arabic{enumi})}
\setcounter{enumi}{13}
\item
  Let \(\mathrm{X} = [0,1]\) be equipped with the Lebesgue measure, and let \(\mathrm{N}\colon \mathrm{X}\times \mathrm{X}\to \mathbf{R}_{+}\) be such that \(\mathrm{N}(x,y) = |x - y|\) . The endomorphism \(u_{\mathrm{N}}\) of \(\mathrm{L}^2 (\mathrm{X})\) is not positive.
\item
  Let X be a locally compact topological space, let \(\mu\) be a positive bounded measure on X with support X, and let \(N \in \mathcal{C}_{b}(X \times X)\) .
\end{enumerate}

\begin{enumerate}
\def\labelenumi{\alph{enumi})}
\tightlist
\item
  There exists an endomorphism \(u\) of \(\mathcal{C}_b(\mathrm{X})\) such that
\end{enumerate}

\[
u (f) (y) = \int_ {\mathrm{X}} \mathrm{N} (x, y) f (x) d \mu (x)
\]

for every function \(f \in \mathcal{C}_b(\mathrm{X})\) and every \(y \in \mathrm{X}\) . In what follows, we assume that \(u\) is compact; this is the case, for example, if the space \(\mathrm{X}\) is compact.

\begin{enumerate}
\def\labelenumi{\alph{enumi})}
\setcounter{enumi}{1}
\item
  The image of \(u_{\mathrm{N}}\) is contained in \(\mathcal{C}_b(\mathrm{X})\) and \(u_{\mathrm{N}}\) defines a continuous linear map from \(\mathrm{L}_{\mathbf{C}}^{2}(\mathrm{X})\) into \(\mathcal{C}_b(\mathrm{X})\) .
\item
  Let \(f \in \mathrm{L}_{\mathbf{C}}^{2}(\mathrm{X})\) and \(\lambda \in \mathbf{C}^*\) such that \(u_{\mathrm{N}}(f) = \lambda f\) . Then \(f \in \mathcal{C}_b(\mathrm{X})\) .
\end{enumerate}

\(d)\) There exists a countable family \((f_i)_{i\in I}\) in \(\mathcal{C}_b(\mathrm{X})\) and a family \((\lambda_i)_{i\in I}\) of nonzero real numbers such that for all \(f\in \mathrm{L}_{\mathbf{C}}^2 (\mathrm{X})\) , we have

\[
u (f) = \sum_ {i \in I} \lambda_ {i} \langle f _ {i} | f \rangle f _ {i}
\]

in \(\mathcal{C}_b(\mathrm{X})\) .

\begin{enumerate}
\def\labelenumi{\alph{enumi})}
\setcounter{enumi}{4}
\tightlist
\item
  Suppose that \(u_{\mathrm{N}}\) is a positive endomorphism of \(\mathrm{L}_{\mathbf{C}}^{2}(\mathrm{X})\) . For every \(i \in \mathrm{I}\) , we denote by \(h_i \in \mathcal{C}_b(\mathrm{X} \times \mathrm{X})\) the function defined by \(h_i(x,y) = \overline{f_i(x)} f_i(y)\) . Then
\end{enumerate}

\[
N = \sum_ {i \in I} \lambda_ {i} h _ {i}
\]

in \(\mathcal{C}(\mathrm{X}\times\mathrm{X})\) endowed with the topology of compact convergence. (Prove that the series converges using Dini's theorem (TG, X, p. 34, th. 1), then prove the equality.)

\begin{enumerate}
\def\labelenumi{\arabic{enumi})}
\setcounter{enumi}{15}
\tightlist
\item
  Let \(\mathrm{X} = [0,1]\) endowed with the Lebesgue measure.
\end{enumerate}

\begin{enumerate}
\def\labelenumi{\alph{enumi})}
\item
  For any sequence \((f_n)_{n\in \mathbf{N}}\) of continuous functions on X that is orthonormal in \(\mathrm{L}_{\mathbf{C}}^2 (\mathrm{X})\) , there exists a kernel \(k\in \mathcal{C}(\mathrm{X}\times \mathrm{X})\) such that
\item
  The endomorphism \(u_{k}\) of \(\mathrm{L}_{\mathbf{C}}^{2}(\mathrm{X})\) is positive;
\item
  For every \(n \in \mathbf{N}\) , the function \(f_n\) is an eigenfunction of \(u_k\) whose eigenvalue is nonzero;
\item
  The family \((f_n)\) is a basis of the image of the endomorphism \(u_k\) .
\item
  There exist kernels \(k \in \mathcal{C}(\mathrm{X} \times \mathrm{X})\) and orthonormal families of eigenfunctions \((f_n)\) of the endomorphism \(u_k\) of \(\mathrm{L}_{\mathbf{C}}^2(\mathrm{X})\) consisting of bounded functions, such that the sequence \((\|f_n\|_\infty)\) is not bounded.
\end{enumerate}

\begin{enumerate}
\def\labelenumi{\arabic{enumi})}
\setcounter{enumi}{16}
\item
  Let E and F be Hilbert spaces. The bilinear map \((u,v)\mapsto\) \(\mathrm{Tr}(uv)\) defines an isometric isomorphism from the Banach space \(\mathcal{L}_1(\mathrm{E};\mathrm{F})'\) into the space \(\mathcal{L}(\mathrm{F};\mathrm{E})\) .
\item
  Let E be a complex Hilbert space and let \(I_{E}\) be the set defined in no. 3 of IV, p. 151.
\end{enumerate}

Let \(n \in I_{E}\) . Let F be a complex Hilbert space. The map \(u \mapsto \lambda_{n}(|u|)\) from \(\mathcal{L}^{c}(E;F)\) into \(R_{+}\) is continuous.

\begin{enumerate}
\def\labelenumi{\arabic{enumi})}
\setcounter{enumi}{18}
\item
  We keep the notations of Lemma 5 of IV, p. 161. Show that if the image of \(f\) is contained in \(\mathbf{R}\) , if \(f\) is strictly convex, and if \(\varrho_n > 0\) , then there is equality in inequality (6) if and only if \(a_i = b_i\) for \(0 \leqslant i \leqslant n\) .
\item
  Let E be a pre-Hilbert space, let F be a Banach space, and let \(u: E \to F\) be a linear map.
\end{enumerate}

\begin{enumerate}
\def\labelenumi{\alph{enumi})}
\item
  Prove that u is a compact linear map if and only if \(\|u(e_{n})\|\to0\) for every orthonormal sequence \((e_{n})\) of E. (To prove that this condition is sufficient, use TG, IX, p. 20, prop. 14.)
\item
  Suppose that E is a Hilbert space of countable type and that F = E. For there to exist an orthonormal basis \((e_{n})_{n\in\mathbf{N}}\) of E such that \(\|u(e_{n})\|\to0\) , it is necessary and sufficient that u not be a Fredholm endomorphism.
\item
  Suppose that E is a Hilbert space and that F = E. If one has \(\langle u(e_{n}) | e_{n} \rangle \to 0\) for every orthonormal sequence \((e_{n})\) of E, then u is a compact linear map.
\end{enumerate}

\begin{enumerate}
\def\labelenumi{\arabic{enumi})}
\setcounter{enumi}{20}
\tightlist
\item
  Let E be a Hilbert space of countable type.
\end{enumerate}

\begin{enumerate}
\def\labelenumi{\alph{enumi})}
\item
  Let A be an infinite set of unit vectors of E which converges weakly to 0 along the filter of complements of finite sets. Prove that there exists a sequence \((x_{n})_{n\geqslant1}\) with values in A and an orthonormal sequence \((e_{n})_{n\geqslant1}\) in E such that \(\|x_{n}-e_{n}\|\leqslant1/n\) for all \(n\geqslant1\) .
\item
  Let u be a linear map from E into E, not necessarily continuous. Suppose that for every orthonormal sequence \((e_{n})\) in E, the sequence \((u(e_{n}))\) converges weakly to 0. Prove that u is continuous and compact.
\item
  Let u be an endomorphism of E. There exists an orthonormal basis \((e_{n})_{n\in\mathbf{N}}\) of E such that the sequence \((u(e_{n}))_{n\in\mathbf{N}}\) converges weakly to 0 if and only if there does not exist an endomorphism v of E such that \(v\circ u-1_{E}\) is compact.
\end{enumerate}

\begin{enumerate}
\def\labelenumi{\arabic{enumi})}
\setcounter{enumi}{21}
\tightlist
\item
  Let E be a complex Hilbert space. Let u and v be positive endomorphisms of E.
\end{enumerate}

\begin{enumerate}
\def\labelenumi{\alph{enumi})}
\tightlist
\item
  For every real number p such that \(0 < p < 1\) and every \(x \in R_{+}\) , we have
\end{enumerate}

\[
x ^ {p} = \frac {\sin (p \pi)}{\pi} \int_ {0} ^ {+ \infty} t ^ {p} \left(\frac {1}{t} - \frac {1}{t + x}\right) d t.
\]

\begin{enumerate}
\def\labelenumi{\alph{enumi})}
\setcounter{enumi}{1}
\item
  Suppose that \(u \leqslant v\) . For every real number \(p\) such that \(0 < p < 1\) , we have \(u^p \leqslant v^p\) . (Use Lemma 14 of I, p. 119, b) and the preceding question.)\\
\item
  Let \(n \in \mathbf{N}\) . We set \(a = v^{n/2} \circ u^n \circ v^{n/2}\) and \(b = (v^{1/2} \circ u \circ v^{1/2})^n\) . If \(a \leqslant 1_{\mathrm{E}}\) , then \(b \leqslant 1_{\mathrm{E}}\) .
\item
  Suppose that u and v have finite trace. For every integer \(n \in N\) , we have
\end{enumerate}

\[
\operatorname{Tr} \left(\left(v ^ {1 / 2} \circ u \circ v ^ {1 / 2}\right) ^ {n}\right) \leqslant \operatorname{Tr} \left(v ^ {n / 2} \circ u ^ {n} \circ v ^ {n / 2}\right).
\]

(Let \(a\) and \(b\) be as in question \(c\) ); by using Lemma 5 of IV, p. 161, reduce to proving that the largest eigenvalue of \(\widehat{\Lambda}^{n}b\) is less than that of \(\widehat{\Lambda}^{n}a\) ; replace E by \(\widehat{\Lambda}^{n}\mathrm{E}\) and \(u\) , \(v\) by \(\widehat{\Lambda}^{n}u\) and \(\widehat{\Lambda}^{n}v\) , respectively, then apply the preceding question.)